\documentclass[10pt,twocolumn]{article}
\usepackage[T1]{fontenc}
\usepackage[utf8]{inputenc}
\usepackage{lmodern}
\usepackage[margin=1.8cm,top=2.4cm,bottom=2cm,columnsep=0.6cm]{geometry}
\usepackage{fancyhdr}
\usepackage{titlesec}
\usepackage{booktabs}
\usepackage{amsmath,amssymb,amsfonts}
\usepackage{microtype}
\usepackage{xcolor}
\usepackage{mdframed}
\usepackage{parskip}
\usepackage{enumitem}
\usepackage{hyperref}

\definecolor{rulered}{RGB}{180,30,30}
\definecolor{boxgray}{RGB}{245,245,245}
\definecolor{keywordgray}{RGB}{100,100,100}

\pagestyle{fancy}
\fancyhf{}
\fancyhead[L]{\textsc{\small Claude Mag}}
\fancyhead[C]{\small\textit{Machine Learning Research Digest}}
\fancyhead[R]{\small 2026-11-25}
\fancyfoot[C]{\small\thepage}
\renewcommand{\headrulewidth}{0.8pt}

\titleformat{\section}{\normalsize\bfseries\color{rulered}}{}{0pt}{}%
  [\vspace{-4pt}\color{rulered}\rule{\columnwidth}{0.4pt}]
\titlespacing{\section}{0pt}{8pt}{4pt}

\hypersetup{colorlinks=true,urlcolor=rulered,linkcolor=black}

\begin{document}
\setlength{\parindent}{0pt}
\setlength{\parskip}{4pt}

{\small\color{keywordgray}\textbf{Keywords:}
  Leaderboard Manipulation \textbullet{} Bradley-Terry \textbullet{} LMArena \textbullet{} Evaluation Integrity \textbullet{} Preference Aggregation}\par\vspace{4pt}

{\LARGE\bfseries\raggedright Leaderboard Hacking: Preference-Based Model Evaluations Are Vulnerable to Manipulation}\par\vspace{4pt}
{\large\itshape\raggedright Three targeted attack strategies can boost a model's arena rank by up to 45 places, exposing Bradley--Terry scoring as structurally gameable}\par\vspace{6pt}
{\small\textbf{NeurIPS 2026 Oral}}\par\vspace{2pt}
{\small\color{keywordgray}arXiv:2611.14769 \textbullet{} cs.LG, cs.AI \textbullet{} NeurIPS 2026 Oral}
\par\vspace{2pt}\noindent{\color{rulered}\rule{\columnwidth}{0.6pt}}\vspace{6pt}

Arena-style leaderboards have become the dominant benchmark for comparing large language models, yet their underlying Bradley--Terry scoring mechanism is structurally vulnerable to three distinct manipulation strategies.
Each strategy alone moves a model 15--21 positions; combined, they produce rank boosts of up to 45 places---enough to elevate a median model into the top five.

\begin{mdframed}[backgroundcolor=boxgray,linecolor=rulered,linewidth=1pt,
    innerleftmargin=6pt,innerrightmargin=6pt,
    innertopmargin=6pt,innerbottommargin=6pt]
\textbf{\small AT A GLANCE}\par\vspace{4pt}
\begin{description}[leftmargin=0pt,labelsep=4pt,itemsep=2pt,parsep=0pt,
    font=\small\bfseries]
  \item[Problem:] Arena leaderboards aggregate pairwise preferences with Bradley--Terry; the statistical mechanism is gameable by any actor who controls even a small fraction of the data.
  \item[Why it matters:] Leaderboard rank is widely used as a proxy for model quality in research, procurement, and media coverage---manipulation distorts all downstream decisions.
  \item[Method:] Three attack strategies: strategic voting, strategic prompting, and strategic nomination; each exploits a different data-injection point.
  \item[Result:] Each attack raises rank by 15--21 places; all three combined raise rank by up to 45 places.
  \item[Evidence:] Simulated LMArena environment; 72 models; 50,000 preference votes; results replicated across starting rank positions.
\end{description}
\end{mdframed}

\section{The Big Picture}

Arena-style leaderboards---most prominently LMArena---have largely displaced fixed benchmark suites as the community's preferred tool for comparing large language models.
Their appeal is intuitive: rather than testing on a frozen dataset that may be contaminated by training data, they aggregate organic human preferences expressed through pairwise battles.
The Bradley--Terry model converts these preferences into a global ranking that is updated continuously as new votes arrive.

This paper argues that the apparent robustness of human-preference aggregation is illusory.
The Bradley--Terry model is a maximum-likelihood estimator trained on the win/loss matrix it is given.
Whoever controls what prompts enter the arena, how votes are cast, and which models compete controls the data the estimator sees.
If that control is exercised strategically, the resulting ranking can be steered far from any honest quality ordering.

\section{Why Existing Approaches Fall Short}

\textbf{Rate limiting and bot detection} constrain the volume of a single attacker's contributions, but the strategies documented here achieve large rank gains with realistic budgets.
A 10\% vote-budget fraction is sufficient to move a model 8--12 places.

\textbf{Rater quality weighting} (downweighting low-quality or inconsistent raters) does not address strategic voting, because an attacker who provides internally consistent but biased votes is indistinguishable from a genuine user with strong preferences.

\textbf{Bradley--Terry extensions} (Davidson ties, weighted variants) inherit the same structural vulnerability because the attack operates on the win-probability matrix, not on specific estimator parameters.

\section{The Three Attack Strategies}

\textbf{Strategic Voting.}
The attacker submits pairwise votes using an oracle policy: upvote the target model against strong competitors (which raises the target's estimated win probability against those models) and vote in unrelated battles in ways that benefit the target's relative standing.
An attacker need not vote in every battle involving the target---targeted votes against the most influential competitors have outsized effect.

\textbf{Strategic Prompting.}
The attacker preferentially submits prompts drawn from a portfolio known to elicit strong responses from the target model.
If the target excels at structured reasoning and poorly-performing models struggle with it, flooding the arena with structured reasoning prompts shifts the battle distribution in the target's favor without casting a single strategic vote.

\textbf{Strategic Nomination.}
The attacker registers decoy models that are systematically weak against specific competitors but strong against the target.
Because Bradley--Terry is a global estimator, the presence of a model that beats competitor $A$ but loses to the target inflates the target's inferred strength relative to $A$---even if the target has never directly competed with $A$.

\section{Technical Formulation}

\subsection*{Bradley--Terry Model}

Let $\beta_i > 0$ be the ``strength'' parameter for model $i$.
The probability that model $i$ beats model $j$ is:
\begin{equation}
P(i \succ j) = \frac{\beta_i}{\beta_i + \beta_j}
\label{eq:bt}
\end{equation}
Given observed win counts $w_{ij}$ (wins of $i$ against $j$), the maximum-likelihood estimates $\hat{\beta}_i$ maximize:
\begin{equation}
\ell(\boldsymbol{\beta}) = \sum_{i,j} w_{ij} \log\frac{\beta_i}{\beta_i + \beta_j}
\label{eq:btll}
\end{equation}

\subsection*{Rank Sensitivity}

The paper shows that a perturbation $\delta W$ to the win matrix produces a rank change that grows as $O(\log N)$ in the number of models $N$.
For an arena of 72 models, a vote budget of 10\% of total votes can achieve a rank change of 8--12 positions through strategic voting alone.

\section{Experimental Findings}

The evaluation uses a simulated LMArena environment: 72 models, 50,000 preference votes, pairwise battles drawn from a Zipf-distributed prompt pool.
The target model starts at median rank (rank 36).

\begin{center}
\small
\begin{tabular}{lrr}
\toprule
\textbf{Attack} & \textbf{Median gain} & \textbf{Max gain} \\
\midrule
Strategic voting & +13.2 & +21 \\
Strategic prompting & +10.8 & +18 \\
Strategic nomination & +7.4 & +15 \\
All three combined & +31.6 & \textbf{+45} \\
\bottomrule
\end{tabular}
\end{center}

Combined attacks are super-additive: the three strategies interact to amplify each other's effects.
Results are robust across starting rank positions from rank~10 to rank~65.

\section{Ablations and Interpretation}

\begin{itemize}[leftmargin=1em,itemsep=2pt]
  \item \textbf{Vote budget:} At 20\% of the optimal vote budget, rank gains are 5--8 positions---well within the range that affects leaderboard visibility.
  \item \textbf{Model pool size:} The logarithmic growth prediction holds empirically: 100-model arenas show 10\% higher maximum rank gains than 72-model arenas.
  \item \textbf{Partial execution:} Running only strategic voting (the simplest attack) achieves 70\% of the combined attack's rank gain.
  \item \textbf{Mitigation attempts:} The authors test three candidate mitigations (vote rate limiting, rater quality weighting, and nomination limits) and find that none reduces rank gains below 5 positions for any single attack.
\end{itemize}

\section{What the Guarantee Says}

There is no formal impossibility result ruling out manipulation-resistant pairwise preference aggregation.
The paper shows that the specific statistical mechanism currently used (Bradley--Terry with maximum-likelihood estimation) is vulnerable, and that the natural mitigations are insufficient.
The authors call for adversarial robustness to be treated as a first-class design criterion alongside statistical efficiency in future leaderboard design.

\section*{Reference}
\textbf{Leaderboard Hacking: Preference-Based Model Evaluations Are Vulnerable to Manipulation.}
arXiv:2611.14769. cs.LG, cs.AI. NeurIPS 2026 (oral).
\url{https://arxiv.org/abs/2611.14769}

\end{document}
